

\section{Racine carrée d'un nombre}

\begin{defi}{}{}
$\sqrt{a}$ est le nombre positif qui a pour carré $a$.
\end{defi}

\medskip
 $a$ est un carré, donc un nombre positif ; 

 ainsi $\sqrt{-5}$ n'existe pas.
\medskip

 Quelques valeurs exactes à connaître:

\vspace{-15mm}
\begin{flushright}
\renewcommand{\arraystretch}{1.5}
\begin{tabular}{|c|c|c|c|c|c|c|c|c|c|c|c|}\hline
$a$&0&1&4&9&16&25&36&49&64&81&100\\\hline
$\sqrt{a}$&0&1&2&3&4&5&6&7&8&9&10\\\hline
\end{tabular}
\end{flushright}

\vspace{-6mm}
\begin{rem}{}{}
Tous les nombres $a$ figurant dans ce tableau sont des carrés parfaits. Plus généralement, on appellera carré parfait les nombres dont la racine carrée est un entier.

\medskip
 Mais une racine carrée ne peut pas toujours s'écrire sous forme d'un entier, d'un décimal ou d'une fraction.

 Par exemple, la seule écriture exacte de la racine carrée de 2 est $\sqrt{2}$. On dit que $\sqrt{2}$ est irrationnel.
\end{rem}

\begin{deme}{}{}
Proposition à démontrer:  $\sqrt{2}$ est irrationnel .

\medskip
On va démontrer cette proposition par l'absurde.

\smallskip
 Supposons que $\sqrt{2}$ soit rationnel. Alors $\sqrt{2}$ est de la forme $\frac{a}{b}$, $a$ et $b$ deux entiers positifs et premiers entre eux. (la fraction $\frac{a}{b}$ ne peut plus être simplifiée)

 On en déduit que: 
$$\sqrt{2} = \dfrac{a}{b}\quad \Longleftrightarrow \quad
\sqrt{2} { \times b}= \dfrac{a}{b}{ \times b}\quad \Longleftrightarrow \quad
\sqrt{2} \times b = a\quad \Longleftrightarrow \quad
\left( \sqrt{2}\times b \right) {^2} =  a {^2}\quad \Longleftrightarrow \quad
\sqrt{2}^2 \times \left( b \right)^2 = a^2\quad \Longleftrightarrow \quad
2 \times b^{2} = a^2 $$
Puisque $a^2$ est pair, $a$ est pair donc $a=2p$ avec $p$ un entier positif
$$2 \times b^{2} = (2p)^2\quad \Longleftrightarrow \quad
2 \times b^{2} = 4p^2\quad \Longleftrightarrow \quad
b^{2} = 2p^2$$

 On en déduit que $b^2$ est pair donc que $b$ est pair.

\medskip
 Par conséquent, il est possible de simplifier la fraction $\frac{a}{b}$ par $2$, ce qui contredit l'hypothèse que $a$, $b$ sont premiers entre eux. 

\medskip
 Puisque l'hypothèse  $\sqrt{2}$ est rationnel  ~~ conduit à une contradiction, c'est le contraire qui est vrai, à savoir: 

 $\sqrt{2}$ est irrationnel .
\end{deme}

\begin{prop}{}{}
Si $a$ et $b$ sont deux entiers naturels, alors $\sqrt{a} \times \sqrt{b} = \sqrt{a\times b}$
\end{prop}

\medskip
\begin{deme}{}{}
Proposition à démontrer: $\sqrt{a} \times \sqrt{b} = \sqrt{a\times b}$

\begin{center}
$\sqrt{a\times b} ^2 \quad = \quad a\times b \quad = \quad \sqrt{a}^2 \times \sqrt{b}^2 \quad = \quad \left( \sqrt{a} \times \sqrt{b} \right) ^2 $
\end{center}

$\sqrt{a\times b}$ \quad et \quad $\sqrt{a} \times \sqrt{b}$ \quad étant positif, on peut en conclure que:
\[\sqrt{a} \times \sqrt{b} = \sqrt{a\times b}\]

\vspace{-8mm}

\end{deme}

\begin{methode}{}{}
Regroupements et simplifications d'écritures:

\begin{itemize}
\item  on cherche à décomposer le nombre qui est sous le radical en un produit dont l'un des facteurs est un carré.
\begin{itemize}
\item[-]$\sqrt{20} = \sqrt{4\times 5} = \sqrt{4} \times \sqrt{5} = 2\sqrt{5}$
\item[-]$\sqrt{32} = \sqrt{16\times 2} = \sqrt{16} \times \sqrt{2} = 4\sqrt{2}$
\item[-]$\sqrt{108} = \sqrt{36\times 3} = \sqrt{36} \times \sqrt{3} = 6\sqrt{3}$
\end{itemize}
\item dans un produit on cherche à associer les racines carrées d'un même nombre:
\begin{itemize}
\item[-]$\sqrt{2}\times \sqrt{14} = \sqrt{2}\times \sqrt{2\times 7} = \sqrt{2}\times \sqrt{2}\times \sqrt{7} = \left(\sqrt{2}\right)^2 \times \sqrt{7} = 2 \sqrt{7}$
\item[-]$\sqrt{3}\times \sqrt{30} \times \sqrt{5}= \sqrt{3}\times \sqrt{3\times 5\times 2} \times \sqrt{5}= \sqrt{3}\times \sqrt{3}\times \sqrt{5} \times \sqrt{5}\times \sqrt{2}= \left(\sqrt{3}\right)^2 \times \left(\sqrt{5}\right)^2 \times \sqrt{2} = 3 \times 5 \times \sqrt{2} = 15\sqrt{2}$
\end{itemize}
\item on cherche à réduire les sommes en mettant les racines identiques en facteur:
\begin{itemize}
\item[-]$3\sqrt{2}+5\sqrt{2} = (3+5)\times \sqrt{2} = 8\sqrt{2}$
\item[-]$3\sqrt{7}-2\sqrt{7} +\sqrt{7}= (3-2+1)\times \sqrt{7} = 2\sqrt{7}$
\end{itemize}
\end{itemize}
Ces calculs ont pour but d'obtenir un résultat dont l'écriture est la plus simple possible.
\end{methode}

%\begin{ex}{}{}
%$\sqrt{3} \times \sqrt{2} = \sqrt{3\times 2} = \sqrt{6} $
%\end{ex}


\begin{rem}{}{}
\begin{itemize}
\item Ce résultat se généralise à tous les nombres positifs et pas uniquement aux entiers.
\item Si de plus $b \ne 0$, alors de même: 

\vspace{-7mm}
\[ \frac{\sqrt{a}}{\sqrt{b}}=\sqrt{\frac{a}{b}}. \]
\item Attention: en général $\sqrt{a} + \sqrt{b}$ N'EST PAS ÉGAL à $\sqrt{a + b}$
\quad($\sqrt{9} + \sqrt{16}= 3 + 4 = 7$ ; mais $\sqrt{9 + 16} = \sqrt{25} = 5$ )
\end{itemize}
\end{rem}


\begin{prop}{}{}[Inégalité triangulaire]
Pour tout nombre strictement positif $a$ et $b$,  $\sqrt{a + b} < \sqrt{a} + \sqrt{b}$
\end{prop}

\begin{deme}{}{}
Proposition à démontrer: $\sqrt{a + b} < \sqrt{a} + \sqrt{b}$
\begin{align*}
\left( \sqrt{a} + \sqrt{b} \right) ^2 \quad = \quad \left( \sqrt{a} + \sqrt{b} \right) \times \left( \sqrt{a} + \sqrt{b} \right) \quad &= \quad \sqrt{a}^2 + \sqrt{a}\sqrt{b} + \sqrt{b}\sqrt{a} + \sqrt{b}^2 \\
&> \sqrt{a}^2 + \sqrt{b}^2 \quad = \quad \quad a+b \quad = \quad \sqrt{a+b} ^2 
\end{align*}
$\sqrt{a+b}$ et $\sqrt{a} + \sqrt{b}$ étant positif, on peut en conclure que:
$\sqrt{a + b} < \sqrt{a} + \sqrt{b}$
\end{deme}

\vspace{-3mm}
	